
\begin{frame}[allowframebreaks]{Definition 2: Equalised Odds (Separation)}

    \begin{block}{Equalised odds --- Hardt, Price and Srebro (2016)}
        \[
            \Pr(\Yhat = 1 \mid Y = y,\, A = \GrpA) \;=\; \Pr(\Yhat = 1 \mid Y = y,\, A = \GrpB)
            \qquad \text{for } y \in \{0,1\}
        \]
        Equal \textbf{TPR} and equal \textbf{FPR} across groups.
    \end{block}

    \vspace{0.5em}

    \begin{block}{Equal opportunity --- the one-sided version}
        Only $\Pr(\Yhat = 1 \mid Y = 1)$ must match: among people who \emph{would} have succeeded,
        each group has the same chance of being selected.
    \end{block}

    \vspace{0.6em}

    \begin{itemize}\small
        \setlength\itemsep{0.35em}
        \item This is the ProPublica criterion: their complaint was unequal FPR and FNR.
        \item Use it when the burden of a \textbf{mistake} is what matters --- a wrongly denied
              loan, a wrongly flagged defendant, a missed diagnosis.
        \item It conditions on $Y$, so it inherits whatever bias is in the label. If $Y$ is
              ``was arrested,'' equalised odds equalises errors against a biased ground truth.
    \end{itemize}

    \vspace{0.3em}
    {\scriptsize M.~Hardt, E.~Price and N.~Srebro, ``Equality of Opportunity in Supervised Learning,''
    NeurIPS 2016, \href{https://arxiv.org/abs/1610.02413v1}{arXiv:1610.02413v1}.}
\end{frame}
